The three other digits are chosen from $\{3,4,5,6,7,8\}$. Since $2+0+1+9=12$, their sum must be divisible by three. The six available digits form residue classes $\{3,6\},\{4,7\},\{5,8\}$, each of size two. A three-element choice with sum zero modulo three must take one from each class, giving $2^3=8$ choices. For each choice, the seven distinct digits have $\dfrac{7!}{4!}=210$ orders with the required relative order of the four fixed digits. Such an order cannot begin with zero because $2$ precedes $0$. Thus the count is $\boxed{8\cdot210=1680}$.
